\section{Analysis and Discussion}
\label{sec:discussion}

\subsection{Why Pre-LayerNorm Matters}
Our experiments confirm the theoretical argument of \textcite{xiong2020layer}:
Pre-LN Transformers train stably at depth without warmup on the residual
path, whereas Post-LN Transformers require careful warmup and often diverge
when the learning rate is set too high. Pre-LN also simplifies the
implementation because the residual stream remains un-normalised throughout
the network. A concrete consequence of this choice is that the loss curve
in Figure~\ref{fig:loss} shows no early-training instability whatsoever:
the very first evaluation point, at iteration 250, already reports a
validation loss of $2.47$, and from there the curve descends monotonically
to $1.545$ without any spikes, plateaus, or regressions. Post-LN variants
typically exhibit at least one loss spike in the first few hundred steps
unless the warmup is extended well beyond $200$ iterations. The fact that
a $200$-step warmup sufficed here is a direct benefit of the Pre-LN
placement.

\subsection{Role of the Causal Mask}
The causal mask is not merely a syntactic constraint---it is what makes
autoregressive generation mathematically consistent with training. During
training, every position predicts the \emph{next} token conditioned on all
previous tokens, computed in a single parallel forward pass thanks to the
mask. This is the key efficiency gain of the Transformer over RNNs, which
must process the sequence sequentially \parencite{vaswani2017attention}.

\noindent We can verify the mask directly by visualising the attention weights.
Figure~\ref{fig:attention} shows the attention matrices of all four heads
in the first Transformer block for the prompt \texttt{"ROMEO: "}. Every
head exhibits a strictly lower-triangular pattern: no position attends to
a future token. Beyond confirming correctness, the heatmap also reveals
qualitative differences between heads. Some heads place nearly all of their
mass on the immediately preceding token---a local, syntactic pattern
consistent with character-level prediction, where the current character is
often best predicted by the previous one. Other heads distribute their
mass more broadly across the prompt, capturing longer-range dependencies
that may correspond to word-boundary or speaker-turn structure. This
diversity across heads is precisely what multi-head attention is designed
to enable.

\subsection{Mixed-Precision Training}
Mixed precision was enabled via \texttt{torch.autocast} with
\texttt{GradScaler}. Two subtleties are worth highlighting. First, under
fp16 the scaler is essential because the small gradient magnitudes from
late-layer parameters underflow fp16's limited exponent range; bf16, which
shares fp32's exponent range, does not require it, so the scaler was
automatically disabled in our run. Second, gradient \emph{clipping} must be
applied \emph{after} unscaling, otherwise the clip threshold is meaningless
because it would be applied to gradients in an inflated numerical range.
Our implementation handles this correctly.

\noindent In terms of throughput, training completed 5000 steps in $3.83$ minutes,
an average of $\approx 21.8$ iterations per second with a steady-state
peak of $\approx 31$~it/s. This rate is dominated by the batch size
($64$) and the model size ($0.82$~M parameters) rather than by arithmetic
precision, which is why a direct fp32 baseline was not measured---at this
scale, the numerical-stability benefit of mixed precision is more
important than raw speed. The end-to-end loss curve shows no NaNs, no
loss spikes, and no evidence of numerical drift, which is the primary
benefit we were seeking.

\subsection{Learning-Rate Schedule}
The combination of linear warmup and cosine decay is empirically the
strongest simple schedule for Transformers \parencite{loshchilov2017sgdr}.
Warmup avoids the early-training instability caused by the random initial
attention patterns; cosine decay allows the model to settle into a
well-conditioned minimum without the sharp drops associated with step
decay.

\noindent The trajectory of validation loss in our run supports this. Between step
$4250$ and step $4750$, the learning rate drops from $4.60 \times 10^{-5}$
to $3.18 \times 10^{-5}$---a factor of roughly $1.45\times$---and
validation loss improves from $1.5590$ to $1.5491$. In the final $250$
steps, the LR reaches its floor of $3.00 \times 10^{-5}$ and validation
loss settles at $1.5453$. This final phase is the smoothest part of the
entire training curve: no oscillation, no regression, and no evidence of
the "bounce" that a poorly-tuned constant learning rate often produces.
The cosine decay did exactly what it is designed to do.

\subsection{Selective Weight Decay}
Theoretically, weight decay is a Gaussian prior on the weights that
encourages small-norm solutions \parencite{loshchilov2019decoupled}.
Applying it to LayerNorm gains is counterproductive because these
parameters are \emph{scale} parameters---pulling them toward zero destroys
the network's ability to renormalise activations. Similarly, decaying bias
terms has no theoretical justification: a bias shift is not a complexity
measure in the way that a weight magnitude is.

\noindent In our optimiser configuration, the parameters were split into
$\mathbf{18}$ tensors with $\dim(\theta) \geq 2$ that received weight decay
$\lambda = 0.1$, and $\mathbf{34}$ one-dimensional tensors---LayerNorm
gains and all biases---that received none. This matches the standard GPT
recipe \parencite{brown2020language}. Because we did not run a control
experiment with weight decay applied to \emph{all} parameters in this
particular training run, we cannot quote a numerical delta in validation
loss as we did in the original proposal. What we \emph{can} say is that
the validation curve remained smooth and monotone throughout, and the
final model achieves a perplexity of $4.69$ on a corpus where a
character-level bigram baseline sits at approximately $3.4$---a
substantial improvement that we attribute in part to this configuration.
Theoretically motivated parameter grouping remains the standard practice
for a reason, and this run is consistent with that consensus.

\subsection{Sampling Trade-offs}
The five generations from Figure~\ref{fig:samples}---all from the same
prompt \texttt{"ROMEO: "}---tell a textbook story about the coherence /
diversity trade-off. Greedy decoding is deterministic and therefore prone
to repetition: once the model enters a high-probability loop, it cannot
escape. In our run, greedy fell into the loop
\emph{``the shall shall be the stand of the stand / The shall be the some
of the some of the son''} within roughly fifteen tokens. This is not a
bug in the model; it is a fundamental property of $\arg\max$ over a
distribution whose modes are locally sticky.

\noindent Temperature scaling is the simplest way to introduce diversity, but it
does so uniformly across the distribution, which also amplifies
low-probability errors. At $T = 0.5$ the output remained locally
grammatical but still repetitive. At $T = 0.8$, combined with top-$k = 40$,
we obtained what we consider the best output of the run: coherent lines,
plausible character names (\emph{``PAULINA''}, \emph{``BUCKINGHAM''}),
and syntactically valid structure, all without obvious looping. This is
what most practitioners consider the sweet spot, and our results agree.

\noindent Pushing temperature higher confirms the failure mode. At $T = 1.0$ the
output begins to introduce non-words and odd morphological variants. At
$T = 1.5$ the text degenerates into something closer to noise than
language: \emph{``'TOnpe shall thou god my fnighthting''}. This is not
because the model is broken but because high temperature flattens the
distribution enough that low-probability continuations---many of which
are simply typos from the training corpus or rare character
combinations---become competitive with plausible ones.

\noindent Top-$k$ sampling sidesteps this by truncating the tail before sampling,
which is why it is standard in most modern language models
\parencite{holtzman2020curious}. It preserves diversity among the
plausible continuations while eliminating the pathological tail. Top-$p$
(nucleus) sampling is a natural extension that adapts $k$ to the shape of
each distribution---keeping more candidates when the model is uncertain,
fewer when it is confident---and would be the next thing we would add.

\subsection{Perplexity in Context}
\label{sec:perplexity}
The final validation perplexity of \textbf{4.69} is worth interpreting
against two reference points. A character-level bigram model on this
corpus achieves roughly $3.4$ perplexity, which is already surprisingly
low because consecutive characters in English have very strong local
correlations: given ``th'', the next character is almost always ``e'',
``a'', or ``o''. A fully-trained character-level LSTM on Shakespeare
typically reaches $\approx 1.8$--$2.0$. Our model sits at $4.69$,
which is worse than the LSTM and---on a per-token basis---only modestly
better than the bigram.

\noindent Why is the improvement over bigram so small relative to the model's
parameter count? Three factors explain this. First, the corpus is
small: $1.1$~M characters is roughly $10\times$ smaller than what LSTM
baselines were tuned on. Second, the model is tiny: $0.82$~M parameters
against $5$--$10$~M for a typical character LSTM. Third, training was
brief: $5000$ steps at batch size $64$ is under $1.3$~M training tokens
seen---roughly $1.2$ epochs over the corpus. The model has simply not
been exposed to the data enough times to fit the higher-order structure
that separates it from a bigram.

\noindent What this result \emph{does} demonstrate, however, is that the
architecture itself is correct. The model captures local structure
(character statistics and short-range dependencies) exactly as expected,
produces coherent short phrases, and generates text that is
unambiguously Shakespeare-shaped rather than random. For a $4$-minute
training run on a laptop, that is the correct operating point. Closing
the gap to the LSTM baseline would primarily require more training
tokens and more epochs, not a change to the architecture.

\subsection{Limitations}
The main limitations of this work are:

\begin{enumerate}
    \item \textbf{Corpus scale.} Tiny Shakespeare is small and
          low-entropy. Absolute perplexity values are therefore not
          directly comparable to those reported on datasets such as
          WikiText-103 or OpenWebText, where character-level models
          routinely reach perplexities in the $1.5$--$2.0$ range.
    \item \textbf{No KV-cache.} Generation currently recomputes the full
          context at every step, giving $O(T^2)$ per generated token
          instead of the $O(T)$ achievable with a KV-cache. For
          $T = 128$ this is imperceptible; for longer contexts it would
          dominate inference time.
    \item \textbf{Hand-tuned hyperparameters.} The learning rate, batch
          size, dropout, and warmup length were chosen by hand based on
          common practice rather than through systematic search. A small
          grid or Bayesian sweep might find a better operating point,
          particularly for the interaction between peak learning rate
          and warmup length.
    \item \textbf{Single seed.} All reported numbers come from a single
          training run. While the loss curve is smooth enough that
          seed-to-seed variance is unlikely to be large, a rigorous
          comparison against a baseline would average over multiple runs.
\end{enumerate}

\noindent None of these limitations undermine the central claim of this report:
that a decoder-only Transformer can be built from scratch, trained
stably with modern optimisation techniques, and produce coherent
autoregressive output---all within a few minutes on commodity hardware.